\section{Conclusion and Future Work}
\label{sec:conclusion}

This project successfully guided a 4-layer Deep Neural Network from a high-level Python algorithm through to a fully verified, tape-out ready Graphic Data System (GDSII) layout in the Nangate 45nm process. By implementing wide 64-bit datapaths, substituting computationally expensive Sigmoids for hardware-friendly ReLU activations, and pushing constant propagation boundaries during logic synthesis, we achieved a highly performant architecture running at an estimated 113.8 MHz with 69.27 mW of power dissipation. Furthermore, tackling the severe routing congestion inherent to dense MLPs by optimizing floorplan density and vertical routing layers resulted in a flawless physical layout with exactly zero DRC violations.

In conclusion, while general-purpose accelerators such as GPUs dominate cloud environments due to their flexible \textit{temporal computing} nature, they remain bottlenecked at the extreme edge due to immense memory bandwidth requirements and instruction decoding latency. GPUs iterate over layers sequentially, fetching varying weights from memory, which allows a single piece of hardware to be multi-purpose and run virtually any model. By deliberately choosing a hardwired, \textit{spatial computing} topology, our ASIC drastically lowered the architectural barrier to entry and memory latency. We traded this multi-purpose operational flexibility for pure, unadulterated efficiency. Our hardware is rigidly fixed to a single 4-layer topology with static weights, but it successfully navigates complex physical layout challenges to yield highly streamlined, ultra-low-latency deep learning silicon specialized exclusively for MNIST.

For future work, the parameterized nature of our SystemVerilog source code allows for effortless architectural scaling. The immediate next step involves performing precise Voltus IR-drop analysis to verify the power grid stability under massive parallel MAC loads, and ultimately isolating the massive weights into an off-chip SRAM wrapper to support significantly larger, multi-megabyte modern AI models.

\subsection{Lessons Learned: Post-Route GLS Debugging}
The transition from RTL verification to Post-Route Gate-Level Simulation (GLS) presented significant challenges that provided valuable insights into the back-end verification flow:
\begin{itemize}
    \item \textbf{Internal Signal Probing:} Physical optimization in Innovus often restructures or deletes internal registers (e.g., FSM state bits). To maintain testbench observability, we learned to identify the actual preserved gate names in the final netlist and use \textit{escaped identifiers} with trailing spaces (e.g., \texttt{u\_dut.\textbackslash state\_reg[0] .Q}) to correctly probe the logic.
    \item \textbf{Simulator Stability Optimization:} Timing violations during the noisy initialization phase can trigger the simulator's \texttt{NOTIFIER} mechanism, propagating `X` states that invalidate the entire run. We stabilized the simulation by utilizing the \texttt{+no\_notifier} flag to suppress premature `X`-propagation and the \texttt{+define+NTC} (Negative Timing Check) macro to allow greater flexibility for post-route SDF delays, ensuring the design reaches a stable functional state.
\end{itemize}

\subsection{Source Code Availability}
The complete project source code, including all RTL, testbenches, synthesis scripts, and APR scripts, is available for review at the following locations:
\begin{itemize}
    \item \textbf{Linux Server:} \texttt{/home/gel8580/NN\_IC\_final\_project}
    \item \textbf{GitHub:} \url{https://github.com/wILLIEWILLYWILLIe/NN_IC_final_project}
\end{itemize}
